% Track 3: user linking (joint). 3c method (C) below; 3b environment signature (A) and 3d hard cases (B)
% to be added by their owners; results to be filled from the campaign (linking/hypotheses.md).

\section{Whose Human? Linking Users}
\label{sec:linking}

\paragraph{Campaign.} Each of us ran the fixed task with our own agents against both partner sites on the
shared VM, as one fictitious persona: every agent three times per site on one network (setting 1, which
also covers the change of site), our own browser on both sites in the same sitting (setting 3), and the
agents again on a second network, a mobile hotspot or a VPN (setting 4). A joint sitting on the campus
network puts all three users behind one address (setting 5). Every visit carries its run identifier in
the path, so the manifest gives the ground-truth user of each request.

\paragraph{Sessions and pairs (3c).} We group requests into sessions: one per run, split into the
agent's requests and any browser requests under the same run identifier. A browser request under an
agent's run identifier can only come from the user opening a link the agent relayed, so such sessions
share the identifier. Each session keeps its client address and its /24 (IPv4) or /64 (IPv6) prefix,
the set of values it showed for each environment-signature field (JA4, JA3, HTTP/2 fingerprint, header
order and set, User-Agent family and version, Accept, Accept-Encoding, key-exchange groups), the
sequence of tools it used, and its timing. For every pair of sessions we compute whether the address and
prefix match, which signature fields share a value, whether the tool sequences match, how close their
pauses are, and whether they share a run identifier.

\paragraph{Linking model.} A logistic regression gives each pair a same-user probability. We evaluate it
with 5-fold cross-validation over sessions in which each pair is scored by a model trained on pairs that
involve neither of its sessions: for folds $k$ and $m$, the pairs between them are scored by a model
trained without both folds. Every pair is then scored out of sample. We report three feature sets:
everything; everything but the address and prefix, which tests linking after an IP change (Track 3 H2);
and the network trace alone, without the address, the relay identifier or the time between sessions,
since that time partly reflects our campaign schedule rather than the user. We group sessions into users
by average-linkage clustering on one minus the pair score and score the clusters with the adjusted Rand
index and B-cubed precision, recall and F1. Connected components (single linkage) are not usable here:
WebFetch sessions of different users look alike, and a single false link merges two users. Baselines are
the same-address rule, the cross-device linking of Zimmeck et al.~\cite{zimmeck2017crossdevice} from
address and browsing history, and re-identification from browsing histories~\cite{bird2020replication}.

\paragraph{Dry run on the local data.} Before the campaign, we ran the pipeline on the 59 local agent
sessions of two users, where every client address is the container's gateway, so it is linking without
the IP. A session that used curl shared its JA4 with every curl session of its user (156 of 156 pairs)
and with none of the other user's (0 of 169), because the system curl carries the machine's TLS stack
(OpenSSL on Linux, Schannel on Windows). WebFetch sessions matched across users as often as within them
(84 of 84 against 87 of 87), and no curl session matched a WebFetch session of the same user. On the
network trace alone the pair model reached an AUC of 0.84. Two users on different operating systems are
the easy case; the campaign results below measure the harder ones.

% TODO (all): campaign results: per setting, linking with and without IP, cross-agent, cross-site,
% agent-browser and IP-change slices; 3b signature stability; 3d shared IP and identical setups.
